\documentclass{article}
\usepackage[T2A]{fontenc}
\usepackage[utf8]{inputenc}   
\usepackage[english, russian]{babel}

% Set page size and margins
% Replace `letterpaper' with`a4paper' for UK/EU standard size
\usepackage[a4paper,top=2cm,bottom=2cm,left=2cm,right=2cm,marginparwidth=1.75cm]{geometry}


\usepackage{amsmath}
\usepackage{graphicx}
\usepackage[colorlinks=true, allcolors=blue]{hyperref}
\usepackage{amsfonts}
\usepackage{amssymb}
% \usepackage[left=1cm,right=1cm,top=1cm,bottom=1cm]{geometry}
\usepackage{hyperref}
\usepackage{seqsplit}
\usepackage[dvipsnames]{xcolor}
\usepackage{enumitem}
\usepackage{algorithm}
\usepackage{algpseudocode}
\usepackage{algorithmicx}
\usepackage{mathalfa}
\usepackage{mathrsfs}
\usepackage{dsfont}
\usepackage{caption,subcaption}
\usepackage{wrapfig}
\usepackage[stable]{footmisc}
\usepackage{indentfirst}
\usepackage{rotating}
\usepackage{pdflscape}
\usepackage{listings}

\usepackage{minted}

\usepackage{MnSymbol,wasysym}

\title{Домашнее задание 6}
\author{Лиза Фоменко}
\date{}

\begin{document}

\maketitle



\section*{\textbf{Задание 1[3]}} \textbf{\textit{На лекции был выведен алгоритм нахождения $a, b$, удовлетворяющих уравнению $ax + by = d$, где $d = \gcd(x, y)$.}}

\begin{enumerate}
    \item \textbf{\textit{Обобщите алгоритм на случай, когда уравнение имеет вид $ax + by = k d$, то есть правая часть не является $\gcd$, а делится на него}}.

Для уравнения вида $ax + by = k d$, где $d = \gcd(x, y)$, мы точно знаем, что левая и правая части делятся на $d$. Мы ничего не можем сходу сказать о делимости левой части на $k$, но мы можем предположить, что $x, y$ обладают свойством делимости на $k$, поэтому от коэффициентов этого требовать не будем. Тогда разделим $ax + by = k d$ на $k$, получим уравнение вида $ax' + by' = d$ - такие уравнения мы умеем решать с помощью расширенного алгоритма Евклида, пусть мы получили частное решение - $x_0, y_0$. Докажем, что, если $x_0, y_0$ - решение уравнения $ax' + by' = d$, то $kx_0, ky_0$ - решение уравнения $ax + by = k d$.

$ax' + by = ax_0 + by_0 = d \Rightarrow k(ax_0 + by_0) = kd = a(kx_0) + b(ky_0) = kd \Rightarrow $ $kx_0, ky_0$ - частное решение исходного уравнения. 


\textbf{тогда алгоритм:} Расширенный алгоритм Евклида для $ax' + by' = d$, полученное частное решение домножаем на $k$ (это будет частное решение исходного уравнения)


    \item \textbf{\textit{Найдите и докажите критерий разрешимости таких уравнений в общем виде. Какие требования нужно наложить на коэффициенты, чтобы уравнение имело решения?}}

Мне кажется на коэффициенты не нужно никаких ограничений, кроме тех, которые накладываются в оригинальном алгоритме Евклида. Для a, b мы хотим уметь считать gcd, значит подойдут любые целые числа, отличные от 0. на $k$ также нет никаких оргичений. Если мы решаем уравнения в целых числах, то $k$ просто должно быть целым. Кажется, что уравнения вида $ax + by = k d$ всегда будут разрешимы.

Другой момент если мы говорим совсем об обобщенном варианте уравнений $ax + by = с$, то здесь критерий следующий: уравнение имеет решение тогда и только тогда, когда $c$ нацело делится на $\gcd(a,b)$. Почему? Левая часть уравнения $ax + by$ нацело делится на $\gcd(a,b)$, так как это наибольший общий делитель $a, b$. Тогда правая часть уравнения тоже должна нацело делиться на $\gcd(a,b)$. 

\textbf{Критерий решения уравнений:} правая часть уравнения нацело делится на $\gcd(a,b)$. 

    \item \textbf{\textit{На лекции и в первых двух пунктах речь шла о частных решениях, рассмотрим также и общий вид решения. Утверждается, что если у линейного диофантова уравнения есть одно решение, у него бесконечно много решений. Например, числа $(-4, 3)$ являются решением уравнения $2 x + 3 y = 1$, но кроме того решениями является пара $(-7, 5)$, пара $(-10, 7)$ и так далее. Выведите общий вид решения линейных диофантовых уравнений.}}

Пусть у нас есть уравнение вида $ax + by = c$ ($c$ нацело делится на $\gcd(a,b)$) и его частное решение $x_0, y_0$. Найдем общее решение уравнения.

$ax_0 + by_0 = c$

Рассмотрим следующую пару: $x = x_0 + \frac{b}{\gcd(a,b)}, y = y_0 - \frac{a}{\gcd(a,b)}$. Подставим в уравнеие:

$$ax + by = c \Rightarrow a*(x_0 + \frac{b}{\gcd(a,b)}) + b*(y_0 - \frac{a}{\gcd(a,b)}) = ax_0 + \frac{ab}{\gcd(a,b)} + by_0 - \frac{ab}{\gcd(a,b)} = ax_0 + by_0 = c$$

то есть новая пара $x, y$ также является решением уравнения.

Теперь рассмотрим  $x = x_0 + k\frac{b}{\gcd(a,b)}, y = y_0 - k\frac{a}{\gcd(a,b)}$, где $k \in \mathbb{Z}$. 

$$ax + by = c \Rightarrow a*(x_0 + k\frac{b}{\gcd(a,b)}) + b*(y_0 - k\frac{a}{\gcd(a,b)}) = ax_0 + k\frac{ab}{\gcd(a,b)} + by_0 - k\frac{ab}{\gcd(a,b)} = ax_0 + by_0 = c$$

То есть общее решение уравнения:
\[
\begin{cases}
x = x_0 + \dfrac{b}{\gcd(a,b)}\,k, \\[6pt]
y = y_0 - \dfrac{a}{\gcd(a,b)}\,k, \\[4pt]
k \in \mathbb{Z}.
\end{cases}
\]

\end{enumerate}


\bigskip

\section*{\textbf{Задание 2[2]}} \textbf{\textit{Решите уравнения в целых числах. Нужно найти все решения, а не только частное.}}

\begin{enumerate}
    \item $238x + 385y = 133$

Сначала найдем gcd(238, 385)
 
 gcd(238, 385) = gcd(238, 385 - 238) =  gcd(238, 147) =  gcd(238 - 147, 147) =  gcd(91, 147) =  gcd(91, 147 - 91) =  gcd(91, 56) =  gcd(35, 56) =  gcd(35, 21) =  gcd(35 - 21, 21) =  gcd(14, 21) =  gcd(14, 21-14) =  gcd(14, 7) =  gcd(7,7) = 7

Проверим, делится ли 133 на 7: 133/7= 19

Тогда мы можем с помощью расширенного алгоритма Евклида решить уравнение $238k + 385l = 7$, а затем домножить обе части уравнения на 19: $238k*19 + 385l*19 = 7*19 = 133$, то есть решением уравнения будут $x = 19k, y=19l$.

Распишем расширенный алгоритм Евклида для уравнения $238k + 385l = 7$

\begin{center}
\begin{tabular}{ |c|c|c|c|c|c| } 
\hline
a & b & $\lfloor \frac{a}{b} \rfloor$ & d & l & k \\
\hline
385 & 238 & 1 & 3 & 13 & -21 \\ 
238 & 147 & 1 & 3 & -8 & 13 \\ 
147 & 91 & 1 & 3 & 5 & -8 \\ 
91 & 56 & 1 & 3 & -3 & 5\\ 
56 & 35 & 1 & 3 & 2 & -3 \\ 
35 & 21& 1 & 3 & -1 & 2 \\ 
21 & 14 & 1 & 3 & 1 & -1 \\ 
14 & 7 & 2 & 3 & 0 & 1 \\ 
7 & 0 & - & 3 & 1 & 0 \\ 
\hline
\end{tabular}
\end{center}

То есть получили, что $385*13 - 238*21 = 7$ (верно!)
Тогда частным решением исходного уравнеия $238x + 385y = 133$ будут $x = 19k, y=19l \Rightarrow x = -399, y = 247$. 

Но нас попросили найти все решения, а я нашла только частное. Общее решение выражается через частное в виде системы:
\[
\begin{cases}
x = x_0 + \dfrac{b}{\gcd(a,b)}\,t, \\[6pt]
y = y_0 - \dfrac{a}{\gcd(a,b)}\,t, \\[4pt]
t \in \mathbb{Z}.
\end{cases}
\]

Тогда, в нашем конкретном случае,

\[
\begin{cases} 
x = -399 + \frac{385}{7} t = -399 + 55t \\ 
y = 247 - \frac{247}{7} t = 247 - 34t \\ 
t \in \mathbb{Z}
\end{cases}
\]


\medskip


    \item $143x + 121y = 52$

Сначала найдем  gcd(143, 121). Не хочу Евклидом, $121 = 11^2 \Rightarrow$ делители 1, 11 и 121. 143 очевидно не делится на 121, но очевидно делится на 11. Поэтому  gcd(143, 121) = 11.

Однако 52 не делится на 11. Это значит, что левая часть уравнения делится на 11 нацело, а правая - нет. Тогда в целых числах не существует решения данного уравнения.
    
\end{enumerate}

\bigskip



\section*{\textbf{Задание 3[1]}} \textbf{\textit{Решите уравнение $68x + 85 \equiv 0 \mod 561$ с помощью расширенного алгоритма Евклида. Требуется найти все решения в вычетах.}}

(С использованием Кормена)

$68x + 85 \equiv 0 \mod 561$

$68x \equiv -85 \mod 561$

$68x \equiv 476 \mod 561$
Это уравнение вида $ax \equiv b \mod n$

Найдем  gcd(a,n) =  gcd(68, 561) = 17. Проверим, что $b = 68$ делится нацело на $d =  gcd(a,n) = 17$ (следствие 31.22). 68/17 = 4, тогда наше уравнение разрешимо. При этом уравнение имеет $d = 17$ различных по модулю (различных по $\mod 561$) решений. 

Решим уравнение вида $ax' + ny' = d$, то есть $68x' + 561y' = 17$, с помощью расширенного алгоритма Евклида


\begin{center}
\begin{tabular}{ |c|c|c|c|c|c| } 
\hline
n & a & $\lfloor \frac{a'}{b'} \rfloor$ & d &y' & x'  \\
\hline
561 & 68 & 8 & 17 & 1 &  -8 \\ 
68 & 17 & 4 & 17 & 0 & 1 \\ 
17 & 0 & - & 17 & 1 & 0 \\ 
\hline
\end{tabular}
\end{center}

Тогда $ax' + ny' = d \Rightarrow 68x' + 561y' = 17$ имеет частное решение $x' = -8, y' = 17$. 

Найдем $x_0$ - частное решение исходного уравнения. $x_0 = x'*(b/d) \mod n = -8 * (476/17) \mod 561 = -8*28 \mod 561 = -224 \mod 561 = 337 \mod 561 = 337$.

Частное решение $x_0 = 337$.

Вспомним, что, зная частное, можно найти все $d = 17$ различных по модулю решений уравнения в виду формулы:
$x_i = x_0 + i\frac{n}{d} = 337 + \frac{561i}{17} = 337 + 33i$, для всех $i = 0, 1, ..., d-1=>16$. 

Я конечно же распиывать это 17 раз не буду) По формуле выше, ответ из 17ти вычетов, каждый из которых является решением уравнения (а также любые числа $x \equiv x_i \mod 561$ для любого $i = 0, ..., 16$):

$$337, 370, 403, 436, 469, 502, 535, 568, 601, 634, 667, 700, 733, 766, 799, 832, 865$$


\bigskip


\section*{\textbf{Задание 4[1]}} \textbf{\textit{Найдите обратный остаток $7^{-1} \mod 102$ с помощью расширенного алгоритма Евклида.}}

ищем $7^{-1} \mod 102$. По определению,  $7^{-1}$ - это такое число, что $7 * 7^{-1} \equiv 1 \mod 102$. Обозначим 
$x = 7^{-1}$ и решим уравнение $7x \equiv 1 \mod 102$.

(Решение с использованием Кормена)

Во-первых, проверим, что gcd(a,n)=gcd(7,102)=1 делит нацело на b=1. Делит, тогда наше уравнение разрешимо, и, согласно Кормену (следствие 31.22) имеет единственное по модулю (!!!) решение. 

Сначала нам необходимо найти решение уравнения $ax' + ny' = gcd(a,n) \Rightarrow 102x' + 7y' = 1$. Решим с помощью расширенного алгоритма Евклида

\begin{center}
\begin{tabular}{ |c|c|c|c|c|c| } 
\hline
n & a & $\lfloor \frac{a'}{b'} \rfloor$ & d &y' & x'  \\
\hline
102 & 7 & 14 & 1 & 2 & -29  \\ 
7 & 4 & 1 & 1 & -1 & 2 \\ 
4 & 3 & 1 & 1 & 1 & -1 \\ 
3 & 1 & 3 & 1 & 0 & 1\\ 
1 & 0 & - & 1 & 1 & 0 \\ 
\hline
\end{tabular}
\end{center}

$102*2 - 7*29 = 1$, тогда $x_0 = x'(b/d) \mod n = -29 * (1/1) \mod 102 = -29 \mod 102 = 2$ является частным решением нашего уравнения, то есть для $7x = 1 \mod 102$ $x_0 = -29$ является частным решением. Приведем к нормальному виду: $-29 \equiv 73 \mod 102$, поэтому $x_0 = 73$ ялвяется частынм решением нашего уравнения. 

Так как у нас $gcd(a,n) = gcd(7,102) = 1$, то решение единственное по модулю, то есть $x = 73$ (тут вопрос о единственности, так как все $x' \equiv 73 \mod 102$ будут решать уравнение. я все таки за единственность обратных объектов))))

\bigskip


\section*{\textbf{Задание 5[2]}} \textbf{\textit{Предложите эффективный алгоритм вычисления наименьшего общего кратного (НОК) двух чисел в битовой модели вычислений (время выполнения операций зависит от длины битовой записи чисел). Докажите его корректность и оцените сложность.}}

Вроде как выводить формулу $lcm(a,b) = \frac{ab}{gcd(a,b)}$ не надо, поэтому будем ею пользоваться как данностью. Еще можно для чисел $a = 2^{k_1} *3^{k_2} * 5^{k_3} ...$ и $b = 2^{l_1} *3^{l_2} * 5^{l_3} ...$ расчитывать $lcm(a,b) = 2^{max(k_1, l_1)} *3^{max(k_2, l_2)} * 5^{max(k_3, l_3)} ...$, после троеточия продолжается последотельность простых чисел $p_i : p_i \leq a//2$.

Наивный алгоритм:

\begin{verbatim}
0 def LCM(a: int (n bit), b: int >= 1 ) -> lcm, int:
1     lcd = Euclide-lcd(a, b)                           # алгоритм Евклида за n^3 (с) лекция
2     ab = Multiply(a, b)                               # функция multiply за n^2 (c) дасгупта
3     lcm = Divide(ab, lcd)                             # функция divide за n^2   (с) предыдущее дз
4     return lcm
\end{verbatim}

Сложность: $O(n^3) + O(n^2) + O(n^2) = O(n^3)$, где $n$ - длина максимального из чисел a, b.

Для корректности работы этого алгоритма мне лишь надо доказать формулу $lcm(a,b) = \frac{ab}{lcd(a,b)}$. Докажем следующим образом: 

$a = 2^{k_1} *3^{k_2} * 5^{k_3} ...$

$b = 2^{l_1} *3^{l_2} * 5^{l_3} ...$

$lcd(a,b) = 2^{min(k_1, l_1)} *3^{min(k_2, l_2)} * 5^{min(k_3, l_3)} ...$

Это разложение на простые множители, и единственное не очевидное тут это почему берем $min$: если мы возьмем хотя бы $min(k_1, l_1)+1$, то одно из чисел не будет делиться на $lcd$, что противоречим условию.

Аналогично для $lcm(a,b) = 2^{max(k_1, l_1)} *3^{max(k_2, l_2)} * 5^{max(k_3, l_3)} ...$. 

тогда из формулы: $lcm = \frac{ab}{lcd} \Rightarrow ab = lcm*lcd \Rightarrow 2^{k_1 + l_1} *3^{k_2 + l_2} * 5^{k_3 + l_3} = 2^{min(k_1, l_1) + max(k_1, l_1)} *3^{min(k_2, l_2) + max(k_2, l_2)} * 5^{min(k_3, l_3) + max(k_3, l_3)} =^{k_1 + l_1} *3^{k_2 + l_2} * 5^{k_3 + l_3}  \Rightarrow$ формула верна, а значит и наш алгоритм тоже. А так как в алгоритме все используемые функции уже были доказаны, то и алгоритм в целом корректен и имеет асимптотику $O(n^3)$

\bigskip




\end{document}